% ── SECTION 1: INTRODUCTION — THE TWO-WINGS TEST (~700 words) ───────────────

\section{Introduction: The Two-Wings Test}
\label{sec:intro}

% Key question: what would genuine confirmation of the two-wings principle
% look like? Not shared conclusions — but shared structure.
% Introduce structural isomorphism vs. analogy.
% Preview six convergences and chronological sequence.

\AbdulBaha taught that science and religion are ``like the two wings upon which
the bird of human intelligence can soar into the heaven of truth.''
The principle is foundational to \bahai teaching and frequently invoked.
But what would it mean to confirm it?

Previous \bahai-science dialogue has often worked by identifying resonances:
quantum uncertainty sounds like the unknowability of God; entanglement sounds
like the interconnection of all things; the Big Bang sounds like creation.
These resonances are suggestive.
But they are analogical, not structural.
Analogy shows that two things look similar; structural isomorphism shows that
they have the same logical form.

This paper asks whether the two-wings principle can be confirmed at the level
of structure --- whether the \bahai metaphysical framework and the framework
that modern physics has empirically arrived at share the same logical
architecture, arrived at by independent means.

The answer I document is yes.
On my reading, the \bahai metaphysics of \mahabbat --- relational love as the creative
foundation of existence at every degree --- shares the same logical structure
as relational quantum mechanics~\citep{Rovelli1996} and process
ontology~\citep{Whitehead1929}: existence constituted by relationship,
properties emerging through relational interaction, the force sustaining
relationship identified as love.

The paper proceeds in five movements.
\S\ref{sec:mahabbat} establishes, on my reading, the \bahai metaphysics of \mahabbat through
close reading of three core passages and their secondary corroborations.
\S\ref{sec:convergences} documents six structural correspondences between that
metaphysics and the framework of contemporary physics: relational ontology,
nonlocal interconnection, observer-dependence and the hierarchy of
comprehension, decoherence and the dissolution of forms, the holographic and
implicate orders, and motion as the character of existence.
\S\ref{sec:chronology} places the convergences on a timeline that begins in
the \bahai revelation of the late nineteenth century, passes through
\citeauthor{Whitehead1929}'s philosophical articulation in 1929, and arrives
at the experimental confirmation that earned the 2022 Nobel Prize in Physics.
\S\ref{sec:confirmation} addresses the methodological question raised by the
convergence and distinguishes it from any claim of borrowing or appropriation.
\S\ref{sec:soul} considers one implication for \bahai scholarship: the
structural status of the non-composite soul within a relational framework.
\S\ref{sec:conclusion} returns to the two-wings principle and what its
confirmation at the level of structure --- as distinct from the level of
metaphor --- contributes to academic \bahai-science dialogue.
